\textbf{1. Kiến thức cần nhớ}
\begin{itemize}
\item \textbf{Hàm số} $y = ax^2$ ($a \ne 0$):
\begin{itemize}
\item Tập xác định: $D = \mathbb{R}$.
\item Tính đơn điệu:
\begin{itemize}
\item Nếu $a > 0$: Hàm số nghịch biến khi $x < 0$, đồng biến khi $x > 0$. Đạt giá trị nhỏ nhất $y = 0$ tại $x = 0$.
\item Nếu $a < 0$: Hàm số đồng biến khi $x < 0$, nghịch biến khi $x > 0$. Đạt giá trị lớn nhất $y = 0$ tại $x = 0$.
\end{itemize}
\end{itemize}
\item \textbf{Đồ thị của hàm số} $y = ax^2$ ($a \ne 0$): Là một đường cong \textbf{Parabol} $(P)$:
\begin{itemize}
\item Có đỉnh là gốc tọa độ $O(0; 0)$.
\item Nhận trục tung $Oy$ làm \textbf{trục đối xứng}.
\item Nếu $a > 0$: Parabol nằm phía trên trục hoành, quay bề lõm lên trên.
\item Nếu $a < 0$: Parabol nằm phía dưới trục hoành, quay bề lõm xuống dưới.
\end{itemize}
\end{itemize}

\textbf{2. Phương pháp giải các dạng}

%% DANG 9D6?18-1
\textbf{Dạng 1. Tính giá trị hàm số, xác định hệ số $a$}
\begin{enumerate}
\item Điểm $M(x_0; y_0) \in (P)\colon y = ax^2 \iff y_0 = a \cdot x_0^2 \implies \boxed{a = \dfrac{y_0}{x_0^2}}$.
\item Biết $x$, thay vào tính $y$; biết $y_0$, giải $ax^2 = y_0 \implies x = \pm\sqrt{\dfrac{y_0}{a}}$ (chú ý có hai giá trị đối nhau).
\end{enumerate}

%% DANG 9D6?18-2
\textbf{Dạng 2. Vẽ đồ thị, đọc đồ thị}
\begin{enumerate}
\item \textbf{Lập bảng giá trị}: Chọn ít nhất $5$ điểm đối xứng nhau qua gốc tọa độ:
$$x \in \{-2k, -k, 0, k, 2k\} \implies y = ax^2.$$
\item Biểu diễn các điểm trên mặt phẳng tọa độ $Oxy$.
\item Vẽ đường cong parabol đối xứng đi qua các điểm đã chọn.
\end{enumerate}

%% DANG 9D6?18-3
\textbf{Dạng 3. Tương giao giữa parabol và đường thẳng}
\begin{enumerate}
\item Cho $(P)\colon y = ax^2$ và đường thẳng $(d)\colon y = mx + n$.
\item \textbf{Phương trình hoành độ giao điểm}:
$$\boxed{ax^2 = mx + n \iff ax^2 - mx - n = 0 \quad (*)}$$
\item Số giao điểm của $(d)$ và $(P)$ bằng số nghiệm của phương trình $(*)$:
\begin{itemize}
\item $\Delta > 0 \iff (d)$ cắt $(P)$ tại hai điểm phân biệt.
\item $\Delta = 0 \iff (d)$ tiếp xúc với $(P)$ tại một điểm duy nhất (tiếp điểm).
\item $\Delta < 0 \iff (d)$ không cắt $(P)$.
\end{itemize}
\end{enumerate}

%% DANG 9D6?18-4
\textbf{Dạng 4. Bài toán thực tế về hàm số $y = ax^2$}
\begin{enumerate}
\item Mô hình hóa các cổng vòm parabol, chuyển động ném ngang, gương cầu, cáp cầu treo\dots
\item Chọn hệ trục tọa độ $Oxy$ thích hợp (thường đặt gốc $O$ tại đỉnh parabol hoặc chân cột đối xứng).
\item Xác định tọa độ điểm đã biết để tìm hệ số $a$, từ đó trả lời câu hỏi thực tế (độ cao, độ rộng, thông số kích thước).
\end{enumerate}
